\subsection{有向集}

定义7. 一个预序集E被称为右有向的（相应地，左有向的），如果E的任意两个元素构成的子集都有上界（相应地，有下界）。

我们常常用“关于关系 \(\leqslant\) 有向”来代替“右有向”，以及当预序关系用其他符号表示时的类似表达式。例如，若S是集合A的子集族，我们说S关于关系 \(\subset\) （相应地 \(\supset\) ）是定向的，如果对于每个由S中两个元素组成的子集 \(\{X, Y\}\) ，存在 \(Z \in S\) 使得 \(X \subset Z\) 并且 \(Y \subset Z\) （相应地 \(X \supset Z\) 并且 \(Y \supset Z\) ）。

示例

\begin{enumerate}
\def\labelenumi{(\arabic{enumi})}
\tightlist
\item
  具有最大元素的有序集是右定向的。
\end{enumerate}

* (2) 在拓扑空间中，一个点的邻域基关于关系 \(\supset\) 是定向的。

\begin{enumerate}
\def\labelenumi{(\arabic{enumi})}
\setcounter{enumi}{2}
\tightlist
\item
  任意模的有限型子模所成的集合关于关系 \(\subset\cdot_{*}\) 是定向的。
\end{enumerate}

命题10. 在右有向有序集E中，一个极大元 \(a\) 是E的最大元。

对于每一个 \(x \in \mathbf{E}\) ，由假设存在 \(y \in \mathbf{E}\) 使得 \(x \leqslant y\) 并且 \(a \leqslant y\) ；由于 \(a\) 是极大元，故 \(y = a\) .

右有向预序集关于相反序是左有向的。右有向集的任何乘积都是右有向的。另一方面，右有向集的子集不一定是右有向的。然而，右有向集E的共尾子集F总是右有向的；对于给定的 \(x, y \in \mathbf{F}\) ，存在 \(z \in \mathbf{E}\) 使得 \(x \leqslant z\) 并且 \(y \leqslant z\) ，然后再取 \(t \in \mathbf{F}\) 使得 \(z \leqslant t\) .

